\documentclass[11pt]{article}
\usepackage[margin=1in]{geometry}
\usepackage[T1]{fontenc}
\usepackage[utf8]{inputenc}
\usepackage{amsmath,amssymb,amsthm}
\usepackage{enumitem}
\usepackage{hyperref}

\newtheorem*{theorem}{Theorem}

\newcommand{\PSh}{\mathsf{PSh}}
\newcommand{\Sh}{\mathsf{Sh}}
\newcommand{\SFin}{\mathsf{Site}}
\newcommand{\Cart}{\mathsf{C}}
\newcommand{\Vcat}{\mathsf{V}}
\newcommand{\pt}{\operatorname{pt}}
\newcommand{\op}{\mathrm{op}}

\title{A linear/non-linear adjunction of sheaves\\ over a site of causal interaction contexts}
\author{Hoapa project}
\date{June 2026 \\[2pt]
\normalsize Extended abstract --- machine-checked in cubical Agda}

\begin{document}
\maketitle

\begin{abstract}
We construct, with a complete machine-checked proof in cubical Agda, a
sheaf-level linear/non-linear adjunction in the sense of Benton over a small
site of causal interaction contexts. The objects of the site are finite local
views of a concurrent computation --- a community of unityped actors, the
events they take part in, and the causal order on those events. The topology is
the one generated by the decomposition of each context into its single events.
We show that a presheaf on this site is a sheaf exactly when it is determined by
its values at single events, and that the inclusion of these sheaves into
presheaves is reflective: for any coefficient category with finite products, the
reflector (sheafification, in the precise sense made explicit below) sends a
presheaf to the product of its event-local values. Given a Benton linear/non-linear
adjunction $F\dashv G$ between a cartesian category and a $*$-autonomous one at the
level of coefficient categories, we lift it to sheaves and obtain an adjunction
between cartesian-valued and linear-valued sheaves, with the linear exponential
``$!$'' recovered as the induced comonad.
\end{abstract}

This note isolates one machine-checked theorem from a larger programme that
models communities of communicating agents --- and the typed data their
interaction produces --- in the language of sheaves and adjoint functors. The
programme rests on a deliberate \emph{type discipline}: the agents themselves are
unityped, indistinguishable to the system, while all structure is carried by the
typed data they produce --- session types in the orchestration layer this theorem
concerns, other type disciplines in other layers. Its
governing commitment is perspectival: there is no absolute,
view-from-nowhere description of what a community of agents is doing, so any
global description must be assembled from local ones. The theorem below is the
first place that commitment is made precise on the operational side. It concerns
a site whose objects are local snapshots of an interaction and a topology whose
only covers are the decompositions of a snapshot into its individual events. In
that setting a \emph{sheaf} is precisely a description determined by what each
single event sees, and \emph{sheafification} is the canonical, universal way to
manufacture such a global-from-local description out of any raw assignment of
meanings to snapshots. The system's ``global point of view'' is thus not imposed
from outside; it is the coherent reconstruction of its event-local standpoints.
The second half of the theorem shows that this reconstruction is compatible with
the resource-sensitive (session-typed) structure the typed data themselves carry.

\section*{The site}

The actor model of concurrency~\cite{Hewitt1973} describes computation as a
community of agents --- \emph{actors} --- that interact solely by exchanging
messages and that the system does not distinguish from one another: they are
\emph{unityped}. To this we add causal structure in the sense of
Lamport~\cite{Lamport1978}: the events of an interaction carry a partial order
recording which event may have influenced which. Following the tradition of
presheaf models of concurrency~\cite{CattaniWinskel1996}, where concurrent
behaviour is studied through (pre)sheaves on a small category of observations, we
organize these local views into a site.

Fix a set $A$ of actors --- a bare set, with nothing distinguishing its elements,
which is the formal content of the actors being \emph{unityped}. A
\emph{finite causal context} over $A$ is a finite set
of events equipped with a partial order --- the causal order --- together with an
incidence relation recording which actors take part in which events; we require
the context to be \emph{supported}, meaning that every actor present takes part
in at least one event. A morphism of contexts is a \emph{causality-reflecting
embedding}: an embedding of one context as a sub-view of another that preserves
actors, events, and incidence, and that both preserves and reflects the causal
order. Write $\SFin_A$ for the resulting category and call it the
\emph{actor-causal site}. This is our name for the object; its ingredients ---
unityped actors, a causal order on events, presheaf semantics --- are standard,
and the causal order is the one familiar from event structures in Winskel's
sense.

The topology is generated by a single kind of cover: the decomposition of a
context into its individual events. For an event $e$ of a context $F$, let
$\pt_e(F)$ be the one-event subcontext at $e$, whose actors are exactly those
that $e$ touches. The generating covering family of $F$ is
\[
  \{\,\pt_e(F)\longrightarrow F\,\}_{e\in E(F)}.
\]
We make no other covers primitive: causal refinements and higher overlaps of
events are deliberately left to richer topologies. We call a presheaf satisfying
descent for these covers a \emph{point sheaf}. We stress at once what is and is
not proved: descent is established for this explicitly presented family of
single-event covers; identifying it with the sheaf condition for a fully
formalized Grothendieck coverage in the sense of SGA\,4~\cite{SGA4} is left to
later work (see Limitations).

\section*{Sheafification}

Let $\mathcal C$ be a \emph{coefficient category} with finite products --- the
category in which our presheaves take their values, and where those values are the
typed data the system assigns to a context. Write
$\PSh_A(\mathcal C)=[\SFin_A^{\op},\mathcal C]$ for the $\mathcal C$-valued
presheaves and $\Sh_A(\mathcal C)\subseteq\PSh_A(\mathcal C)$ for the full
subcategory of point sheaves. A presheaf is a point sheaf precisely when it is
determined by its values at single events; the formalization proves this as an
equivalence between descent for the canonical cover and a product-locality
condition.

The inclusion $\iota_{\mathcal C}:\Sh_A(\mathcal C)\hookrightarrow\PSh_A(\mathcal C)$
has an explicit left adjoint. For a presheaf $P$, set
\[
  a_{\mathcal C}(P)(F)=\prod_{e\in E(F)} P\bigl(\pt_e(F)\bigr),
\]
with restrictions induced by the point-to-point morphisms. This is the
\emph{sheafification} functor for the point topology --- the reflector
exhibiting point sheaves as a reflective subcategory of presheaves. We use the
word ``sheafification'' in exactly this restricted sense, and no other:
$a_{\mathcal C}$ is the left adjoint to the inclusion of the presheaves that
satisfy descent for the canonical single-event covers. We do not claim it is the
associated-sheaf functor for a fully formalized coverage. The construction is
carried out behind an abstract finite-product interface, so that the proof
assistant does not unfold the products into enumeration-dependent normal forms
during unification. The machine-checked result is the adjunction
\[
  a_{\mathcal C}\dashv \iota_{\mathcal C}.
\]

\section*{The linear/non-linear adjunction}

The typed data produced by interaction are resource-sensitive: a message once
consumed is not silently reused, a session once advanced does not revert. This is
the domain of \emph{linear logic}, and --- because the word is badly overloaded
--- we mean linear \emph{logic} throughout, never linear algebra. We take the
linear coefficient category $\Vcat$ to be a $*$-autonomous
category~\cite{Barr1979}, a model of multiplicative linear logic and the
categorical home of session types in the sense of Caires and
Pfenning~\cite{CairesPfenning2010}. The contrasting ``non-linear'' side is the
cartesian fragment, in which typed data may be freely copied and discarded.

A \emph{linear/non-linear adjunction} in the sense of Benton~\cite{Benton1995}
consists of a cartesian category $\Cart$, a symmetric monoidal (here
$*$-autonomous) category $\Vcat$, and a monoidal adjunction
$F\dashv G:\Cart\rightleftarrows\Vcat$; the induced comonad $F\circ G$ on $\Vcat$
is the linear exponential ``$!$'', the modality that re-admits the structural
rules of copying and discarding. We assume such an adjunction at the level of
coefficient categories, with both $\Cart$ and $\Vcat$ carrying finite products.
The formalized hypothesis is in fact a thin, Benton-style signature: it records
the two functors, the adjunction $F\dashv G$, and the strong-monoidal comparison
maps of $F$, while the monoidal coherence and the comonad laws of ``$!$'' are
left unpacked --- the proof consumes only the adjunction. The theorem is thus
more general than a full Benton model would make it.

Postcomposition lifts $F$ and $G$ to presheaves,
\[
  F_* : \PSh_A(\Cart)\to\PSh_A(\Vcat),
  \qquad
  G_* : \PSh_A(\Vcat)\to\PSh_A(\Cart),
\]
with $F_*\dashv G_*$ computed pointwise. The right adjoint preserves point
sheaves: descent along the canonical cover is transported across the natural
bijection $\Cart(X,GY)\cong\Vcat(FX,Y)$. The left adjoint is made to land in
sheaves by composing with sheafification. Writing $R=G_*$ restricted to point
sheaves and $L=a_{\Vcat}\circ F_*\circ\iota_{\Cart}$, we obtain the following.

\begin{theorem}
Let $\Cart$ be a cartesian category and $\Vcat$ a $*$-autonomous category, both
with finite products, and let $F\dashv G:\Cart\rightleftarrows\Vcat$ be a
linear/non-linear adjunction in the sense of Benton. Then there is an adjunction
\[
  L \dashv R : \Sh_A(\Cart)\rightleftarrows\Sh_A(\Vcat)
\]
between cartesian-valued and linear-valued point sheaves on the actor-causal
site. The induced comonad $L\circ R$ is the sheaf-level linear exponential
``$!$''.
\end{theorem}

\begin{sloppypar}
The result is formalized without postulates in the modules
\texttt{SupportedChangeOfBase}, \texttt{SupportedLinearNonlinearSheafAdapter},
and \texttt{PointLinearNonlinearAdjunction}, building on the reflector
constructed in \texttt{PointSheafAdjunction}.
\end{sloppypar}

\section*{Scope and significance}

The theorem is deliberately modest, and three classical lineages meet in it. It
is a descent statement over a small site, in the tradition of
Grothendieck~\cite{SGA4,MacLaneMoerdijk1992}. It is an instance of Lawvere's
principle that adjointness is the right way to organize foundational
structure~\cite{Lawvere1969}: the global object is the value of a left adjoint,
and its universal property is exactly the freeness of the reconstruction. And it
lifts Benton's linear/non-linear adjunction from coefficients to sheaves. The
conceptual content is that event-local observations \emph{freely} generate a
coherent global object, and that this free reconstruction lifts the
linear/non-linear structure from values to sheaves. In the wider programme, this fixes
the operational side of a conjectured larger adjunction between concurrent
practice and modal foundations; the present theorem makes that operational side
concrete, with descent, sheafification, and the linear/non-linear structure all
machine-checked.

\section*{Limitations}

We are explicit about what is \emph{not} done here.
\begin{itemize}[leftmargin=1.4em,itemsep=2pt]
\item The value-level adjunction $F\dashv G$ is \emph{assumed}, not constructed.
\item No concrete session-type model is yet exhibited for $\Vcat$ together with a
verification that it has the required finite products; the theorem is generic in
its hypotheses.
\item The topology is the event-discrete one: causal order and overlaps of events
do not yet carry descent content. Enriching the site so that they do is the
natural next step, and is where colimit and exactness questions return.
\item Point sheaves are defined by descent for the explicitly presented
single-event covers; we do not yet prove that they coincide with the sheaves for
a fully formalized Grothendieck coverage, nor compare against a library coverage.
\item This is not the cross-site adjunction to the modal/foundational side; it is
the first checked result on the operational side alone.
\end{itemize}

\section*{A note on authorship}

This abstract was drafted with the assistance of a language model and reviewed by
the authors; the surrounding framework and its philosophical motivation are
presented separately. The mathematics it reports, by contrast, is not
provisional: every claim above is backed by a proof that the cubical Agda
type-checker accepts, without recourse to postulates in the relevant modules.

\begin{thebibliography}{9}

\bibitem{Hewitt1973}
C. Hewitt, P. Bishop, and R. Steiger.
A universal modular ACTOR formalism for artificial intelligence.
In \emph{Proceedings of the 3rd International Joint Conference on Artificial
Intelligence (IJCAI)}, pp. 235--245, 1973.

\bibitem{Lamport1978}
L. Lamport.
Time, clocks, and the ordering of events in a distributed system.
\emph{Communications of the ACM} 21(7):558--565, 1978.

\bibitem{CattaniWinskel1996}
G. L. Cattani and G. Winskel.
Presheaf models for concurrency.
BRICS Report Series RS-96-35, 1996. Also in
\emph{Computer Science Logic (CSL '96)}, LNCS 1258, Springer, 1997.

\bibitem{SGA4}
M. Artin, A. Grothendieck, and J.-L. Verdier, eds.
\emph{Th\'eorie des topos et cohomologie \'etale des sch\'emas (SGA 4)}.
Lecture Notes in Mathematics 269, 270, 305. Springer, 1972--1973.

\bibitem{MacLaneMoerdijk1992}
S. Mac Lane and I. Moerdijk.
\emph{Sheaves in Geometry and Logic: A First Introduction to Topos Theory}.
Springer, 1992.

\bibitem{Lawvere1969}
F. W. Lawvere.
Adjointness in foundations.
\emph{Dialectica} 23(3--4):281--296, 1969.

\bibitem{Benton1995}
P. N. Benton.
A mixed linear and non-linear logic: proofs, terms and models.
In \emph{Computer Science Logic (CSL '94)}, LNCS 933, pp. 121--135. Springer,
1995.

\bibitem{Barr1979}
M. Barr.
\emph{$*$-Autonomous Categories}.
Lecture Notes in Mathematics 752. Springer, 1979.

\bibitem{CairesPfenning2010}
L. Caires and F. Pfenning.
Session types as intuitionistic linear propositions.
In \emph{CONCUR 2010}, LNCS 6269, pp. 222--236. Springer, 2010.

\end{thebibliography}

\end{document}
